\section{Experimental method}
\label{sec:method}
\subsection{Machine, configurations, and version separation}
Table~\ref{tab:environment} identifies the physical machine and the main current
configuration. The physical CPU has heterogeneous cores; the reported guest
affinities are virtual processor identities, not a claim of fixed placement on
particular physical core types. Firmware menu settings were not directly
photographed or queried. Windows feature fields that report virtualization as
false beneath an active hypervisor are retained literally and are not interpreted
as proof that firmware VMX is disabled.

\begin{table}[tbp]
\centering\small
\caption{Single-machine environment. The Hyper-V compatibility experiment uses
a separate Windows~1 boot from the VMware and performance experiments.}
\label{tab:environment}
\begin{tabularx}{\linewidth}{@{}lX@{}}
\toprule
Component & Recorded configuration \\
\midrule
Physical CPU & Intel Core i7-13700F; family 6, model 183, stepping 1;
16 cores, 24 logical processors \\
Outer Windows & Windows 11 Pro for Workstations Insider, 26300.9022 \\
Outer Hyper-V & \code{hvix64.exe} 10.0.26100.9022; outer HVCI running \\
Firmware & ASUS S501ME.331; exact firmware VMX/VT-d menu state unobserved \\
Windows 1 & Windows 11 Pro 22621.4317; 4 vCPUs, fixed 8 GiB;
Generation 2 outer VM, nested extensions exposed \\
VMware & Workstation 17.6.4, build 24832109; virtual hardware version 21 \\
Descendant & TinyCore 17.1; Linux 6.18.35-tinycore64; normal \code{/init};
2 online vCPUs, 768 MiB; reconfigured to 4096 MiB for the reachability and
region sweeps, which state which \\
Inner Hyper-V fixture & Generation 1, 2 vCPUs, fixed 768 MiB;
diskless TinyCore evidence ISO \\
\bottomrule
\end{tabularx}
\end{table}

We separate four evidence scopes. The \emph{current Pro cohort} uses the
implementation subsequently committed as \code{bd703d84}, the version~3 lease,
and the driver hash beginning \code{3b44132f}; it supplies every performance,
latency and lifecycle number. The \emph{region cohort} is a later build of the
same driver carrying the version~4 lease, and supplies only the large-leaf,
reachability, staging, digest and drift results; no performance number is
carried across from it, and none of the current cohort's numbers were re-measured
on it. The \emph{historical 4-by-2 cohort}
uses Windows~1 Home before the Pro upgrade, implementation \code{b9bdda83}, and
driver hash beginning \code{8850e28f}. It contains a broader benchmark suite,
high-load page tests, and owner-exit tests. The inner Hyper-V admission experiment
uses the preceding ordinarily test-signed driver, whose exact identity is
recorded with that experiment. Its refusal precedes resource preparation.
Full hashes and the evidence revision appear in Appendix~\ref{app:artifact}.
The parent commit of a build made from a modified working tree is not treated
as the complete measured source identity; provenance also records input hashes.

Changing Windows edition, processor count, or inner Hyper-V startup required
target reboots during preparation. These reboots are outside measured continuity
intervals. The outer Hyper-V and HVCI remain active. For VMX-dependent experiments,
inner Hyper-V startup is disabled before a separate boot; it is restored after
cleanup. No experiment silently pools these boots or treats setup reboots as
live transitions.

The driver was built with the standard MSVC/WDK toolchain and separately signed
with the existing ordinary laboratory test certificate. Variant signing was
disabled. Driver loading on the target succeeded; an untrusted certificate root
reported by host Authenticode verification is not represented as production
trust certification. The current Pro cohort's driver was not rebuilt for
preparation of this manuscript. The region cohort's was: its results come from
binaries built and deployed during this work, and Appendix~\ref{app:artifact}
names each by the hash of the file the service actually loaded --- because one
of them turned out not to behave like a clean rebuild of the same sources, which
is reported in Section~\ref{sec:limitations} rather than smoothed over here.

\subsection{Continuity, clocks, and observation}
The harness records Windows boot time, process PIDs and creation times, and
per-CPU resident identities around insertion/removal. Descendant page runs add
the VMware VMX process identity, guest boot ID and uptime, and CPU-pinned
observers. HTTP runs additionally record the HTTP server PID/start ticks and
fresh guest pagemap observations. These are sampled continuity checks over
explicit intervals. They are stronger than process-name equality, but they are
not a cryptographic proof that no unobserved state reconstruction could occur.

QPC supplies Windows control, driver, and application timestamps. The analysis
checks that the driver boundaries fall inside the control-call window. Guest
uptime timestamps and serial output provide guest-side ordering; we do not
subtract timestamps from different clock domains to infer sub-millisecond
causal delays. TSC-cycle profiles use a separate invariant-TSC calibration.
CPU rendezvous envelopes and entry-return skew describe internal synchronization,
not a measured universal application stop time.

\subsection{Workloads and repetition}
The current performance experiment has six modes and two workloads: CPUID
leaf~0 latency and TCP loopback RTT. Each mode/workload has one warmup and five
measured repetitions, for 72 runs. Modes are deterministically shuffled within
rounds using seed 20260916. All modes share a driver, controller, workload binary,
and Windows boot, with descendant VMMs absent. The sparse/full nested comparison
preserves nested policy and CPUID visibility; the artificial 256/512-VMREAD modes
use ordinary residency and estimate the marginal cost of extra reads plus loop
overhead.

The separate disruption experiment uses a paced, closed-loop TCP client on
Windows vCPU~2, an echo server on vCPU~1, and control on vCPU~0. It includes four
conditions: sham query while off, insertion, sham query while on, and removal.
Each condition has one warmup and five measured runs. Every response records
sequence, request begin/end QPC values, and the previous completion. Pacing uses
\code{Sleep(1)}. A command-overlap completion gap is a response-completion interval
intersecting the control call; its duration includes pacing and scheduling as
well as any effect of insertion.

Page lifecycle trials require both TinyCore CPUs to read the complete 4096-byte
page and match the expected MD5, together with complete event traces and
allocation accounting. MD5 is used as a convenient non-adversarial equality
check, not a collision-resistant integrity claim. The current cohort repeats
each of five request-local fault modes three times and retains each run
individually. Hardware invalidation failures are not manufactured by these
software fault modes.

The historical nested benchmark comparison uses the same static workload binary
and pins it to guest CPU~0, with one warmup and three measured repetitions per
workload. Its driver versions and guest boots differ. The historical Windows
comparison uses off/pre, on, off/post blocks with seven measured runs per block
after warmup. Reported bootstrap intervals describe resampling within that
machine and boot; they are not independent-machine confidence bounds. We report
medians and observed ranges for the smaller current samples without claiming
statistical significance.

\subsection{Retention and exclusions}
Every run is retained as machine-readable JSON/JSONL or raw serial plus a
per-run record. Warmups remain in the archive but are excluded from reported
performance medians. Incomplete HTTP observation, setup failures, malformed
requests, and rejected admission remain visible in their own categories.
An earlier benchmark batch overlapping a host build is excluded from the
historical main comparison and retained. For the guest direct-write application
comparator, the analysis excludes six observations using the stated 0.5-second
stage settling rule; Section~\ref{sec:application} explains its limits. There is
no aggregate ``success rate'' combining different configurations, software
versions, setup attempts, and observation contracts.
